\documentclass[12pt,a4paper]{article}
\usepackage{amsmath,amssymb,amsthm,hyperref,geometry}
\usepackage{enumitem,booktabs,array}
\usepackage[T1]{fontenc}
\usepackage[utf8]{inputenc}
\usepackage[ngerman]{babel}
\geometry{margin=2.5cm}

% -------------------------------------------------------
% Theorem-Umgebungen
% -------------------------------------------------------
\newtheorem{theorem}{Satz}[section]
\newtheorem{lemma}[theorem]{Lemma}
\newtheorem{corollary}[theorem]{Korollar}
\newtheorem{proposition}[theorem]{Proposition}
\newtheorem{conjecture}[theorem]{Vermutung}
\newtheorem{definition}[theorem]{Definition}

\theoremstyle{remark}
\newtheorem{remark}[theorem]{Bemerkung}
\newtheorem{example}[theorem]{Beispiel}

% -------------------------------------------------------
% Metadaten
% -------------------------------------------------------
\title{\textbf{Die Birch-und-Swinnerton-Dyer-Vermutung}\\[0.5em]
\large Rang, $L$-Funktionen und die Arithmetik elliptischer Kurven}
\author{Michael Fuhrmann}
\date{März 2026}

\hypersetup{
  pdftitle={Die Birch-und-Swinnerton-Dyer-Vermutung},
  pdfauthor={Michael Fuhrmann},
  colorlinks=true,
  linkcolor=blue,
  citecolor=blue
}

\begin{document}
\maketitle
\begin{abstract}
Die Vermutung von Birch und Swinnerton-Dyer (BSD), in den frühen 1960er Jahren auf Basis
umfangreicher numerischer Experimente formuliert, sagt eine tiefe Verbindung zwischen dem
algebraischen Rang einer elliptischen Kurve über $\mathbb{Q}$ und der Nullstellenordnung
ihrer zugehörigen $L$-Funktion bei $s=1$ voraus. Sie ist eines der sieben
Millennium-Preisaufgaben des Clay Mathematics Institute.

Dieses Paper stellt den Satz von Mordell--Weil vor, der die Gruppenstruktur
$E(\mathbb{Q}) \cong \mathbb{Z}^r \oplus T$ begründet, die genaue BSD-Vermutung in
schwacher und starker Form, die bekannten Teilergebnisse (Rang 0 via Coates--Wiles 1977,
Rang 1 via Gross--Zagier 1986 und Kolyvagin 1989), numerische Verifikationen aus der
Cremona-Datenbank und dem LMFDB, die Verbindung zur Iwasawa-Theorie sowie den aktuellen
Stand des Problems im Jahr 2026.
\end{abstract}

\tableofcontents
\newpage

% ================================================================
\section{Elliptische Kurven und die Mordell--Weil-Gruppe}
% ================================================================

\subsection{Grundlegende Definitionen}

\begin{definition}[Elliptische Kurve]
  Eine \emph{elliptische Kurve} über $\mathbb{Q}$ ist eine glatte projektive Kurve vom
  Geschlecht $1$ mit einem ausgezeichneten rationalen Punkt $O$ (dem Neutralelement).
  Sie lässt sich in \emph{Weierstraß-Form} schreiben:
  \begin{equation}\label{eq:weierstrass}
    E\colon y^2 + a_1 xy + a_3 y = x^3 + a_2 x^2 + a_4 x + a_6,
    \quad a_i \in \mathbb{Z},
  \end{equation}
  oder in \emph{kurzer Weierstraß-Form} (nach quadratischer und kubischer Ergänzung):
  \[
    y^2 = x^3 + Ax + B, \quad A,B\in\mathbb{Q},\; \Delta := -16(4A^3 + 27B^2) \neq 0.
  \]
\end{definition}

\begin{definition}[Gruppengesetz]
  Die rationalen Punkte $E(\mathbb{Q})$ bilden eine abelsche Gruppe gemäß dem
  Sekanten-Tangenten-Verfahren: Durch zwei Punkte $P, Q$ auf $E$ legt man eine Gerade,
  die $E$ in einem dritten Punkt $R$ trifft; die Summe ist $P + Q := -R$ (Spiegelung an
  der $x$-Achse). Das Neutralelement ist der Punkt im Unendlichen $O$.
\end{definition}

\subsection{Der Satz von Mordell--Weil}

\begin{theorem}[Mordell--Weil~\cite{Mordell1922,Weil1928}]
  Sei $E$ eine elliptische Kurve über $\mathbb{Q}$. Dann ist die Gruppe der rationalen
  Punkte $E(\mathbb{Q})$ eine endlich erzeugte abelsche Gruppe:
  \begin{equation}\label{eq:mordellweil}
    E(\mathbb{Q}) \;\cong\; \mathbb{Z}^r \;\oplus\; T,
  \end{equation}
  wobei $r \geq 0$ der \emph{Rang} von $E$ und $T$ eine endliche Gruppe, die sogenannte
  \emph{Torsionsuntergruppe}, ist.
\end{theorem}

\begin{theorem}[Mazurs Torsionssatz~\cite{Mazur1977}]
  Die Torsionsuntergruppe $T = E(\mathbb{Q})_{\mathrm{tors}}$ ist isomorph zu einer
  der folgenden fünfzehn Gruppen:
  \[
    \mathbb{Z}/n\mathbb{Z} \text{ für } n \in \{1,2,3,\ldots,10,12\},
    \quad\text{oder}\quad
    \mathbb{Z}/2\mathbb{Z} \times \mathbb{Z}/2n\mathbb{Z} \text{ für } n \in \{1,2,3,4\}.
  \]
\end{theorem}

Der Rang $r$ misst die „Größe" von $E(\mathbb{Q})$: Für $r = 0$ gibt es nur endlich viele
rationale Punkte; für $r \geq 1$ gibt es unendlich viele.

\subsection{Berechnung des Rangs: Abstieg}

Der Standardalgorithmus für den Rang verwendet den \emph{$2$-Abstieg}:
\begin{enumerate}[nosep]
  \item Berechne die $2$-Selmer-Gruppe $\mathrm{Sel}^{(2)}(E/\mathbb{Q})$.
  \item Der Selmer-Rang liefert eine obere Schranke:
    $r \leq \dim_{\mathbb{F}_2}\mathrm{Sel}^{(2)} - \dim_{\mathbb{F}_2} E(\mathbb{Q})[2]$.
  \item Die Shafarevich--Tate-Gruppe $\Sha(E/\mathbb{Q})$ misst das Hindernis:
  \[
    0 \to E(\mathbb{Q})/2E(\mathbb{Q}) \to \mathrm{Sel}^{(2)}(E) \to \Sha(E)[2] \to 0.
  \]
\end{enumerate}

Die BSD-Vermutung sagt voraus, dass $\Sha(E/\mathbb{Q})$ endlich ist; dies ist nur in
Spezialfällen bekannt.

% ================================================================
\section{Die $L$-Funktion einer elliptischen Kurve}
% ================================================================

\subsection{Lokale Faktoren}

Für jede Primzahl $p$ definiert man den lokalen Faktor:
\[
  L_p(E, s) \;=\;
  \begin{cases}
    (1 - a_p p^{-s} + p^{1-2s})^{-1} & \text{falls } E \text{ gute Reduktion bei } p \text{ hat,} \\
    (1 - p^{-s})^{-1}                 & \text{falls } E \text{ gespaltene mult.\ Reduktion hat,} \\
    (1 + p^{-s})^{-1}                 & \text{falls } E \text{ nicht-gespaltene mult.\ Red.\ hat,} \\
    1                                  & \text{falls } E \text{ additive Reduktion hat,}
  \end{cases}
\]
wobei die \emph{Frobenius-Spur} $a_p = p + 1 - \#E(\mathbb{F}_p)$ die Hasse-Schranke
$|a_p| \leq 2\sqrt{p}$ erfüllt.

\subsection{Die globale $L$-Funktion}

\begin{definition}[$L$-Funktion von $E$]
  Die \emph{$L$-Funktion} von $E/\mathbb{Q}$ ist das Euler-Produkt
  \begin{equation}\label{eq:lfunction}
    L(E,s) \;=\; \prod_p L_p(E,s),
  \end{equation}
  das für $\operatorname{Re}(s) > 3/2$ konvergiert.
\end{definition}

Durch den Modularitätssatz (Wiles 1995, Breuil--Conrad--Diamond--Taylor 2001) ist jede
elliptische Kurve über $\mathbb{Q}$ modular: $L(E,s)$ stimmt mit der $L$-Funktion einer
gewichts-$2$-Spitzenform $f_E$ des Niveaus $N$ (dem Führer von $E$) überein:
\[
  L(E,s) \;=\; L(f_E,s) \;=\; \sum_{n=1}^\infty \frac{a_n}{n^s}.
\]
Dies erlaubt die analytische Fortsetzung auf ganz $\mathbb{C}$ und begründet die
Funktionalgleichung:
\begin{equation}\label{eq:functional}
  \Lambda(E,s) \;:=\; \left(\frac{\sqrt{N}}{2\pi}\right)^s \Gamma(s) L(E,s)
  \;=\; \varepsilon(E) \cdot \Lambda(E, 2-s),
\end{equation}
wobei $\varepsilon(E) \in \{+1,-1\}$ die \emph{Wurzelzahl} (root number) von $E$ ist.

\begin{remark}[Paritätsvermutung]
  Die Paritätsvermutung sagt $(-1)^r = \varepsilon(E)$ voraus. Dies gilt unter BSD und
  ist in vielen Fällen unbedingt bewiesen.
\end{remark}

% ================================================================
\section{Die Birch-und-Swinnerton-Dyer-Vermutung}
% ================================================================

\subsection{Historischer Ursprung}

In den frühen 1960er Jahren berechneten Birch und Swinnerton-Dyer~\cite{BSD1965} das
Produkt $\prod_{p \leq X} \#E(\mathbb{F}_p)/p$ für viele Kurven und beobachteten empirisch,
dass dieses Produkt wie $(\log X)^r$ für $X \to \infty$ wächst. Diese Heuristik übersetzt
sich in die folgenden Vermutungen.

\subsection{Schwache BSD-Vermutung}

\begin{conjecture}[Schwache BSD~\cite{BSD1965,BSD1963}]
  Sei $E/\mathbb{Q}$ eine elliptische Kurve vom Rang $r$. Dann gilt:
  \begin{equation}\label{eq:weakbsd}
    \operatorname{ord}_{s=1} L(E,s) \;=\; r \;=\; \operatorname{rang}(E(\mathbb{Q})).
  \end{equation}
  Mit anderen Worten: der analytische Rang stimmt mit dem algebraischen Rang überein.
\end{conjecture}

\subsection{Starke BSD-Vermutung}

Die starke (vollständige) BSD-Vermutung spezifiziert zusätzlich den Leitkoeffizienten
der Taylor-Entwicklung $L(E,s) = c_r (s-1)^r + c_{r+1}(s-1)^{r+1} + \ldots$ bei $s=1$:

\begin{conjecture}[Starke BSD~\cite{BSD1965,Tate1966}]
  \begin{equation}\label{eq:strongbsd}
    \lim_{s\to 1} \frac{L(E,s)}{(s-1)^r}
    \;=\;
    \frac{|\Sha(E)| \cdot \Omega_E \cdot R_E \cdot \prod_p c_p}
         {|E(\mathbb{Q})_{\mathrm{tors}}|^2},
  \end{equation}
  wobei:
  \begin{itemize}[nosep]
    \item $|\Sha(E)|$ die Ordnung der Shafarevich--Tate-Gruppe (vorhergesagt endlich) ist,
    \item $\Omega_E = \int_{E(\mathbb{R})} |\omega|$ die reelle Periode
          (Integral des Néron-Differentials) ist,
    \item $R_E = \det\bigl(\langle P_i, P_j \rangle\bigr)_{1\leq i,j\leq r}$ der
          Regulator (Determinante der Höhen-Paarungsmatrix auf freien Erzeugern) ist,
    \item $c_p = [E(\mathbb{Q}_p) : E_0(\mathbb{Q}_p)]$ die Tamagawa-Zahlen
          (endlich, gleich $1$ für $p \nmid N$) sind,
    \item $|E(\mathbb{Q})_{\mathrm{tors}}|$ die Ordnung der Torsionsuntergruppe ist.
  \end{itemize}
\end{conjecture}

Die Formel \eqref{eq:strongbsd} sagt gleichzeitig voraus:
\begin{itemize}[nosep]
  \item die Nullstellenordnung von $L(E,s)$ bei $s=1$,
  \item die Endlichkeit von $\Sha(E)$,
  \item den genauen Leitkoeffizienten in Termen arithmetischer Invarianten.
\end{itemize}

% ================================================================
\section{Bekannte Ergebnisse}
% ================================================================

\subsection{Rang 0: Coates--Wiles (1977)}

\begin{theorem}[Coates--Wiles~\cite{CoatesWiles1977}]
  Sei $E/\mathbb{Q}$ eine elliptische Kurve mit komplexer Multiplikation (CM) durch den
  Ganzheitsring eines imaginär-quadratischen Körpers $K$. Falls $L(E,1) \neq 0$, so ist
  $E(\mathbb{Q})$ endlich (d.\,h.\ $r = 0$).
\end{theorem}

Dies war das erste unbedingte Ergebnis in Richtung BSD und verwendet die Theorie der
elliptischen Einheiten und die Iwasawa-Theorie für CM-Körper.

\subsection{Rang 0: Kolyvaginsches Euler-System (1989)}

\begin{theorem}[Kolyvagin~\cite{Kolyvagin1989}]
  Sei $E/\mathbb{Q}$ eine modulare elliptische Kurve (heute für alle $E/\mathbb{Q}$
  bekannt). Falls $L(E,1) \neq 0$, gilt:
  \begin{enumerate}[nosep]
    \item $r = 0$ (endliche Mordell--Weil-Gruppe), und
    \item $\Sha(E/\mathbb{Q})$ ist endlich.
  \end{enumerate}
\end{theorem}

Kolyvaginsches Methode führt \emph{Euler-Systeme von Heegner-Punkten} ein und schränkt
die Selmer-Gruppe durch explizite kohomologische Konstruktionen ein.

\subsection{Rang 1: Gross--Zagier und Kolyvagin (1986--1989)}

\begin{theorem}[Gross--Zagier~\cite{GrossZagier1986}]
  Falls $L(E,1) = 0$, aber $L'(E,1) \neq 0$ (analytischer Rang $1$), hat der
  Heegner-Punkt $y_K \in E(K)$ (für einen Hilfskörper $K$ imaginär-quadratisch) unendliche
  Ordnung.
\end{theorem}

\begin{corollary}[Gross--Zagier + Kolyvagin~\cite{Kolyvagin1989}]
  Unter denselben Voraussetzungen ($\operatorname{ord}_{s=1} L(E,s) = 1$) gilt:
  \begin{enumerate}[nosep]
    \item $r = 1$,
    \item $\Sha(E/\mathbb{Q})$ ist endlich.
  \end{enumerate}
\end{corollary}

Zusammen beweisen diese Ergebnisse: \emph{Analytischer Rang $\leq 1$ $\Rightarrow$ algebraischer Rang $=$ analytischer Rang}.

\subsection{Teilergebnisse für höhere Ränge}

Für analytischen Rang $\geq 2$ ist die Situation schwieriger. Zhang~\cite{Zhang2001} und
andere haben die Gross--Zagier-Formel auf Shimura-Kurven verallgemeinert. Skinner--Urban
~\cite{SkinnerUrban2014} bewiesen eine Teilbarkeit der BSD-Formel mittels Iwasawa-Theorie
für die $p$-adische $L$-Funktion in vielen Fällen.

% ================================================================
\section{Numerische Evidenz}
% ================================================================

\subsection{Cremona-Datenbank und LMFDB}

Die \emph{Cremona-Datenbank}~\cite{Cremona2023} listet elliptische Kurven über $\mathbb{Q}$
nach Führer geordnet auf und enthält:
\begin{itemize}[nosep]
  \item algebraischen Rang (berechnet via $2$- und $4$-Abstieg),
  \item $L(E,1)$ auf hohe Genauigkeit (via Buhler--Gross-Numerik),
  \item Tamagawa-Zahlen, Perioden und Regulator.
\end{itemize}

Die \emph{LMFDB} ($L$-Functions and Modular Forms Database,~\cite{LMFDB}) enthält
stand 2026 Daten zu über $3{,}5$ Millionen elliptischen Kurven.

\subsection{Beispiele numerischer Verifikation}

\begin{table}[h]
\centering
\begin{tabular}{lllll}
\toprule
Kurve (Label) & $N$ & $r$ & $L^{(r)}(E,1)/r!$ & $|\Sha|$ (vorhergesagt) \\
\midrule
\texttt{37a1}     & 37     & 1 & $0{,}3059997\ldots$ & 1 \\
\texttt{389a1}    & 389    & 2 & $0{,}1598058\ldots$ & 1 \\
\texttt{5077a1}   & 5077   & 3 & $1{,}7327\ldots$    & 1 \\
\texttt{234446a1} & 234446 & 4 & berechnet           & 1 (vorhergesagt) \\
\bottomrule
\end{tabular}
\caption{Elliptische Kurven mit verifizierten BSD-Daten (Cremona/LMFDB)}
\end{table}

Für die Kurve \texttt{37a1} ($y^2 + y = x^3 - x$) gilt $L(E,1)=0$,
$L'(E,1) \approx 0{,}306$. Die Heegner-Punkt-Methode konstruiert explizit einen rationalen
Punkt unendlicher Ordnung, was $r=1$ und die starke BSD-Formel bestätigt.

% ================================================================
\section{Iwasawa-Theorie und $p$-adische $L$-Funktionen}
% ================================================================

\subsection{Überblick über die Iwasawa-Theorie}

Die Iwasawa-Theorie untersucht die Variation arithmetischer Objekte in
$\mathbb{Z}_p$-Erweiterungen. Für eine elliptische Kurve $E/\mathbb{Q}$ und eine
Primzahl $p$ betrachtet man die zyklotomische $\mathbb{Z}_p$-Erweiterung
$\mathbb{Q}_\infty = \bigcup_n \mathbb{Q}(\zeta_{p^n})$.

\begin{definition}[$p$-adische $L$-Funktion]
  Die $p$-adische $L$-Funktion $L_p(E, s)$ (Mazur--Swinnerton-Dyer,~\cite{MazurSD1974})
  ist eine $p$-adisch analytische Funktion, die Spezialwerte $L(E, \chi, 1)$ für
  Dirichlet-Charaktere $\chi$ von $p$-Potenz-Führer interpoliert.
\end{definition}

\subsection{Die Iwasawa-Hauptvermutung für elliptische Kurven}

\begin{conjecture}[Iwasawa-Hauptvermutung für $E$]
  Sei $X_\infty$ der Pontryagin-Dual der $p^\infty$-Selmer-Gruppe von $E$ über
  $\mathbb{Q}_\infty$. Dann gilt als Moduln über der Iwasawa-Algebra $\Lambda = \mathbb{Z}_p[[T]]$:
  \[
    \mathrm{char}_\Lambda(X_\infty) \;=\; (L_p(E, s))
    \quad \text{in } \Lambda \otimes \mathbb{Q}_p.
  \]
\end{conjecture}

\begin{theorem}[Skinner--Urban~\cite{SkinnerUrban2014}]
  Für $p \geq 5$, $E$ ohne CM, $p \nmid N \cdot |E(\mathbb{Q})_{\mathrm{tors}}|^2 \cdot \prod_\ell c_\ell$
  und unter der Voraussetzung, dass die $p$-adische $L$-Funktion nicht identisch null ist,
  gilt eine Teilbarkeit der Iwasawa-Hauptvermutung:
  \[
    (L_p(E, s)) \;\mid\; \mathrm{char}_\Lambda(X_\infty).
  \]
\end{theorem}

\subsection{Verbindung zu BSD}

Die Iwasawa-Hauptvermutung impliziert BSD im Rang-0-Fall (wenn $L(E,1) \neq 0$) für gute
gewöhnliche Primzahlen $p$, da die $p$-adische $L$-Funktion dann bei $s=1$ nicht
verschwindet. Für Kurven vom Rang $\geq 1$ verschwindet $L_p(E,s)$ bei $s=1$, und die
Verbindung zu BSD erfordert die Theorie der $p$-adischen Höhen und des $p$-adischen
Regulators.

% ================================================================
\section{Schwache versus starke BSD}
% ================================================================

\subsection{Was die schwache BSD besagt}

Die schwache BSD (Gleichung \eqref{eq:weakbsd}) ist eine reine Aussage über den
analytischen Rang: $L(E,s)$ verschwindet genau in Ordnung $r$ bei $s=1$. Dies ist
für jeden $r \geq 2$ noch offen.

\subsection{Was die starke BSD hinzufügt}

Die starke BSD (Gleichung \eqref{eq:strongbsd}) sagt zusätzlich den \emph{genauen Wert}
des Leitkoeffizienten voraus, der folgendes kodiert:
\begin{itemize}[nosep]
  \item die Endlichkeit und die genaue Ordnung von $\Sha(E/\mathbb{Q})$,
  \item den kanonischen Höhen-Paarungsregulator,
  \item alle Tamagawa-Zahlen und Perioden.
\end{itemize}

Die numerische Verifikation der starken BSD-Formel erfordert:
\begin{enumerate}[nosep]
  \item Berechne $L^{(r)}(E,1)/r!$ mit ausreichender Genauigkeit (via modulare Symbole).
  \item Berechne $\Omega_E$, $R_E$, $\prod c_p$, $|E(\mathbb{Q})_{\mathrm{tors}}|$.
  \item Leite $|\Sha(E)|$ ab und verifiziere, dass es ein Quadrat ist (wie durch die
        Cassels--Tate-Paarung vorhergesagt).
\end{enumerate}

\subsection{Die Shafarevich--Tate-Gruppe}

$\Sha(E/\mathbb{Q})$ klassifiziert \emph{prinzipale homogene Räume} (Torsoren) von $E$,
die lokal trivial sind (Punkte über allen $\mathbb{Q}_v$ haben), aber keinen rationalen Punkt
besitzen:
\[
  \Sha(E/\mathbb{Q}) \;:=\; \ker\!\left(H^1(\mathbb{Q}, E) \to \prod_v H^1(\mathbb{Q}_v, E)\right).
\]
Die Cassels--Tate-Paarung macht $\Sha$ zu einer alternierenden, nicht-entarteten Gruppe,
was $|\Sha|$ (falls endlich) zu einem Quadrat zwingt.

% ================================================================
\section{Aktueller Stand (2026)}
% ================================================================

\subsection{Was bekannt ist}

\begin{itemize}[nosep]
  \item \textbf{Analytischer Rang $0$ $\Rightarrow$ algebraischer Rang $0$:} Bewiesen für
    alle modularen Kurven (= alle $E/\mathbb{Q}$) durch Kolyvagin (1989).
  \item \textbf{Analytischer Rang $1$ $\Rightarrow$ algebraischer Rang $1$:} Bewiesen durch
    Gross--Zagier + Kolyvagin (1986--1989).
  \item \textbf{Endlichkeit von $\Sha$:} In den Rängen $0$ und $1$ bekannt (Kolyvagin).
  \item \textbf{Eine Teilbarkeit (Iwasawa):} Skinner--Urban (2014) für viele Primzahlen $p$.
  \item \textbf{$p$-Umkehrsätze:} Wei Zhang (2014)~\cite{WeiZhang2014} bewies die
    $p$-Umkehrung: algebraischer Rang $0$ $\Rightarrow$ $L(E,1) \neq 0$ unter milden
    Voraussetzungen.
\end{itemize}

\subsection{Was offen bleibt}

\begin{itemize}[nosep]
  \item \textbf{Analytischer Rang $\geq 2$:} Kein unbedingtes Ergebnis, das analytischen
    und algebraischen Rang für $r \geq 2$ verbindet.
  \item \textbf{Starke BSD für Rang $\geq 2$:} Die genaue Formel für den Leitkoeffizienten
    ist unbewiesen.
  \item \textbf{Endlichkeit von $\Sha$ allgemein:} Offen für $r \geq 2$.
  \item \textbf{Unbeschränkter Rang:} Ob $r$ beliebig groß sein kann, ist offen
    (Elkies hat Kurven vom Rang $\geq 28$).
  \item \textbf{Effektive Höhenschranken:} Kein effektiver Algorithmus zur Bestimmung aller
    rationalen Punkte auf einer Kurve vom Rang $\geq 2$ ist bekannt.
\end{itemize}

\subsection{Rechnerische Verifikationen (2026)}

Stand 2026 wurde die BSD für alle elliptischen Kurven in der Cremona-Datenbank
($N \leq 500\,000$) und für Millionen von Kurven im LMFDB numerisch mit hoher Genauigkeit
verifiziert. In allen Fällen:
\begin{itemize}[nosep]
  \item Stimmt der analytische Rang mit dem durch Abstieg berechneten Rang überein.
  \item Gilt die starke BSD-Formel bis zur berechneten Genauigkeit.
  \item Ist $|\Sha|$ in allen Fällen, in denen es berechnet werden kann, ein Quadrat.
\end{itemize}

Kein Gegenbeispiel zu BSD wurde je gefunden.

% ================================================================
\section{Verbindungen zu anderen Millennium-Problemen}
% ================================================================

\begin{itemize}[nosep]
  \item \textbf{Riemann-Hypothese:} Die Hasse-Schranke $|a_p| \leq 2\sqrt{p}$ ist eine
    Konsequenz der Riemann-Hypothese für $L$-Funktionen von Kurven über endlichen Körpern
    (Weil 1948), einem bewiesenen Analogon der RH.
  \item \textbf{Langlands-Programm:} BSD fügt sich in die allgemeine Langlands-Philosophie
    ein, dass $L$-Funktionen, die Motiven zugeordnet sind, Arithmetik kodieren. Der
    Modularitätssatz, der $L(E,s)$ mit Modulformen verbindet, ist ein Meilenstein im
    Langlands-Programm.
  \item \textbf{Hodge-Vermutung:} BSD für abelsche Varietäten über Zahlkörpern ist mit der
    Hodge-Vermutung für abelsche Varietäten verbunden (via der Tate-Vermutung).
\end{itemize}

% ================================================================
\section{Schlussfolgerung}
% ================================================================

Die BSD-Vermutung verkörpert eines der tiefsten Rätsel der Arithmetik: Warum sollte
eine analytische Invariante (die Nullstellenordnung von $L(E,s)$) einer rein algebraischen
Größe (dem Rang der Mordell--Weil-Gruppe) entsprechen? Die Teilergebnisse von
Coates--Wiles, Gross--Zagier und Kolyvagin haben die Tür geöffnet, aber die vollständige
Vermutung — insbesondere für Rang $\geq 2$ — bleibt unbewiesen.

Die präzise Formel der starken BSD verknüpft Perioden, Höhen, Tamagawa-Zahlen und die
rätselhafte Shafarevich--Tate-Gruppe zu einer einzigen Identität und legt tiefgreifende
Verbindungen zwischen Analysis, Geometrie und Algebra nahe, die erst ansatzweise
verstanden werden.

\begin{thebibliography}{99}

\bibitem{BSD1963}
B.~J.~Birch und H.~P.~F.~Swinnerton-Dyer,
\textit{Notes on elliptic curves I},
J.\ reine angew.\ Math.\ \textbf{212} (1963), 7--25.

\bibitem{BSD1965}
B.~J.~Birch und H.~P.~F.~Swinnerton-Dyer,
\textit{Notes on elliptic curves II},
J.\ reine angew.\ Math.\ \textbf{218} (1965), 79--108.

\bibitem{Tate1966}
J.~Tate,
\textit{On the conjecture of Birch and Swinnerton-Dyer and a geometric analog},
Séminaire Bourbaki, Exposé \textbf{306} (1966), 415--440.

\bibitem{Mordell1922}
L.~J.~Mordell,
\textit{On the rational solutions of the indeterminate equations of the third and fourth degrees},
Proc.\ Cambridge Philos.\ Soc.\ \textbf{21} (1922), 179--192.

\bibitem{Weil1928}
A.~Weil,
\textit{L'arithmétique sur les courbes algébriques},
Acta Math.\ \textbf{52} (1928), 281--315.

\bibitem{Mazur1977}
B.~Mazur,
\textit{Modular curves and the Eisenstein ideal},
Publ.\ Math.\ IHES \textbf{47} (1977), 33--186.

\bibitem{CoatesWiles1977}
J.~Coates und A.~Wiles,
\textit{On the conjecture of Birch and Swinnerton-Dyer},
Invent.\ Math.\ \textbf{39} (1977), 223--251.

\bibitem{GrossZagier1986}
B.~H.~Gross und D.~B.~Zagier,
\textit{Heegner points and derivatives of $L$-series},
Invent.\ Math.\ \textbf{84} (1986), 225--320.

\bibitem{Kolyvagin1989}
V.~A.~Kolyvagin,
\textit{Euler systems},
in \textit{The Grothendieck Festschrift}, Bd.~II, Birkhäuser, 1989, S.~435--483.

\bibitem{MazurSD1974}
B.~Mazur und P.~Swinnerton-Dyer,
\textit{Arithmetic of Weil curves},
Invent.\ Math.\ \textbf{25} (1974), 1--61.

\bibitem{SkinnerUrban2014}
C.~Skinner und E.~Urban,
\textit{The Iwasawa main conjecture for $\mathrm{GL}_2$},
Invent.\ Math.\ \textbf{195} (2014), 1--277.

\bibitem{WeiZhang2014}
W.~Zhang,
\textit{Selmer groups and the indivisibility of Heegner points},
Camb.\ J.\ Math.\ \textbf{2} (2014), 191--253.

\bibitem{Zhang2001}
S.~Zhang,
\textit{Gross--Zagier formula for $\mathrm{GL}(2)$},
Asian J.\ Math.\ \textbf{5} (2001), 183--290.

\bibitem{Cremona2023}
J.~Cremona,
\textit{Algorithms for Modular Elliptic Curves},
Cambridge University Press, 2.~Aufl., 1997; Datenbank laufend aktualisiert,
\url{https://johncremona.github.io/ecdata/}.

\bibitem{LMFDB}
The LMFDB Collaboration,
\textit{The $L$-Functions and Modular Forms Database},
\url{https://www.lmfdb.org}, abgerufen 2026.

\bibitem{Silverman2009}
J.~H.~Silverman,
\textit{The Arithmetic of Elliptic Curves},
2.~Aufl., Springer GTM \textbf{106}, 2009.

\bibitem{Wiles1995}
A.~Wiles,
\textit{Modular elliptic curves and Fermat's Last Theorem},
Ann.\ of Math.\ \textbf{141} (1995), 443--551.

\end{thebibliography}

\end{document}
